% =====================================================================
%  CH. 5 LA SENTINELLE · CH. 6 LES SECRETS ET LEURS GARDES
% =====================================================================
\section{La sentinelle — du démon à la serrure}\label{sec:sentinelle}

\subsection{Pourquoi ce mot revient, et ce qu'il a voulu dire trois fois}

« Sentinelle » a désigné trois choses différentes dans l'histoire du coffre, et la
confusion des trois est la question la plus fréquente du terrain. Une fois pour
toutes :

\begin{table}[htbp]
\centering\small
\caption{Les trois vies du mot « sentinelle ».}
\label{tab:sentinelles}
\begin{tabularx}{\textwidth}{@{}l Y Y@{}}
\toprule
\textbf{Époque} & \textbf{Ce que c'était} & \textbf{Sa limite} \\
\midrule
CF2, démon résident (D86) & un processus userland de longue durée : surveille la
couche C, croise deux seuils (processus non blanchi + lecture de contenu), brûle les
fragments & usurpable par un faux heart-beat ; tuable par l'adversaire avant
destruction ; exige un moteur de plateforme (inotify/\code{/proc} au laboratoire,
SACL/ETW attendus sur Windows) \\
CF2, PORTE D120 & le verrou qui refuse de déverrouiller sans sentinelle vivante — la
demande du terrain : « si elle n'est pas lancée, tous nos efforts sont inutiles » &
sur plateforme sans moteur, la porte \emph{avertit} à chaque déverrouillage (« SANS
sentinelle ») : l'ère intermédiaire est honnête mais incomplète \\
CF3, sentinelle instantanée (D121) & \textbf{par invocation} : les cinq vérifications
ordonnées courent DANS chaque acte qui déverrouille — ouvrir, clôture,
\code{verifier}, re-lecture à froid & aucune : elle ne dépend d'aucun état de machine
falsifiable — \emph{le piège est le témoin, et il est hors machine} \\
\bottomrule
\end{tabularx}
\end{table}

La question du terrain, posée à la passation : \emph{« on avait pas décidé de la
supprimer, en migrant en CF3 ? Si on impose CF3 dès la création de projet, on n'a pas
besoin de sentinelle, non ? »} La réponse tient en trois temps. Oui, le \textbf{démon}
est mort en CF3 — remplacé, pas retiré. L'exigence D120 (« indispensable au
déverrouillage ») \emph{reste}, réalisée par invocation : la sentinelle instantanée
n'est pas un composant qu'on peut oublier de lancer, \emph{c'est la serrure elle-même}
— le sine qua non est satisfait par construction, sans démon ni moteur de plateforme à
attendre. Et oui, \textbf{un coffre né CF3 n'a jamais connu de sentinelle à lancer} :
\code{init} le crée à l'époque 1 avec son média du piège, et chaque acte porte sa
propre vérification.

\subsection{La sentinelle instantanée — l'ordre est la loi}\label{sec:sentinst}

À tout acte qui déverrouille, cinq vérifications courent \emph{dans cet ordre} —
l'ordre n'est pas esthétique, il sépare les modes de défaillance :

\begin{enumerate}
\item \textbf{Piège présent.} Absent : refus nommé, déverrouillage \emph{impossible}
(fail-closed). Rien n'est lu, rien ne brûle — l'absence du média n'est jamais une
attaque, c'est un état.
\item \textbf{Piège intègre} (le \code{tag}). Invalide : refus nommé + journal, mais
\emph{pas} de destruction — le média est hors limite de confiance : une corruption
bénigne (lecteur capricieux, secteur mort) doit pouvoir être relue depuis les copies
hors-ligne. La destruction attend la preuve cryptographique (étape 5).
\item \textbf{Cohérence des champs publics.} Époque du piège = époque du conteneur
(sinon : « piège PÉRIMÉ », refus anti-retour) ; H\_projet du piège = H\_projet du
conteneur (sinon : le piège d'un \emph{autre} coffre).
\item \textbf{Phrase} (via \code{verif}). Divergente : réessai à délai croissant —
1, 2, 4 secondes, plafond 8 (le plafond borne aussi le verrouillage de compte par
soi-même). \emph{Jamais} de destruction : une faute de frappe ne doit pas tuer le
coffre.
\item \textbf{AEAD + étiquette} (KEK complet). Échec \emph{après} 1--4 verts :
falsification ou corruption \emph{prouvée} côté machine $\to$ destruction locale +
PV + journal. L'étiquette vérifiée (le clair doit hacher vers le H\_projet de
l'en-tête) ferme la course de clôture.
\end{enumerate}

Un échec \textbf{nomme son étape} ; les étapes suivantes ne courent pas. C'est la
grammaire même du fail-closed : vous saurez toujours \emph{quoi} a refusé, et donc
\emph{où} regarder (chapitre 9).

\begin{terrain}[— l'avertissement D120 dans un rapport réel]
Pendant la migration du coffre réel (30/10), le \emph{déverrouillage CF2} — la
dernière fois que la vieille serrure parlait — a affiché :
\begin{lstlisting}
[coffre-CF2] AVERTISSEMENT D120 : deverrouillage SANS
sentinelle vivante - le moteur n'est pas livre sur cette
plateforme (exigence consignee spec  9) : le coffre s'ouvre
sans temoin.
\end{lstlisting}
Aucun avertissement équivalent n'est apparu aux actes CF3 : la sentinelle instantanée
n'a pas de moteur à attendre. L'avertissement D120 est le signal d'une \emph{ère} —
honnête, bornée, et close par la génération suivante.
\end{terrain}

\subsection{Ce que la destruction locale ne touche JAMAIS}

La réponse à la falsification (étape 5) détruit \emph{côté machine} : le conteneur, le
staging, le brouillon — trois passes d'écrasement (aléa, 0xFF, 0x00), troncature,
renommage aléatoire, unlink, fsync, et un PV horodaté dans \code{.keys/pv-cf3/} avec la
raison nommée (deux destructions la même seconde produisent deux fichiers distincts —
correctif D123, appris d'un écrasement prouvé par le terrain). Elle ne touche
\textbf{jamais} le média du piège (c'est la relève), \textbf{jamais} les fragments CF2
ni leurs copies hors-ligne : la destruction reste survivable — c'est la contrepartie
qui rend le piège destructif acceptable. La mort d'un coffre est toujours un état
\emph{recouvrable par le sachet}.

\section{Les secrets et leurs gardes}\label{sec:gards}

\subsection{L'inventaire honnête — ce que vous protégez, en fait}

\begin{table}[htbp]
\centering\small
\caption{Les secrets du coffre et leurs gardes — l'inventaire complet.}
\label{tab:inventaire}
\begin{tabularx}{\textwidth}{@{}l Y Y@{}}
\toprule
\textbf{Secret} & \textbf{Sa garde} & \textbf{Si la garde manque} \\
\midrule
La phrase de contrôle & votre tête ; le séquestre D183/D223 (cinq parts
Shamir à seuil trois, la reconstitution à froid) & sans la phrase, pas de
K\_user — le coffre est muet ; sous trois parts assemblables, la phrase est
PERDUE (le seuil est la survie) \\
Les phrases par projet & le coffre chiffré (CF1/2/3) & un projet dont la phrase est
perdue doit faire l'objet d'une rotation de sa propre lignée \\
Le secret maître M (CF2) & cinq fragments Shamir, seuil 3 — jamais entier nulle part &
sous 3 fragments : mort \emph{informationnelle} (aucune attaque possible, aucune
reconstruction non plus) \\
S\_piège (CF3, par époque) & le média du sachet — jamais sur la machine hors l'instant
de l'acte & sans le piège d'époque égale : l'acte est \emph{impossible} — branchez le
bon média \\
La graine de clé I & support physique hors-ligne (64 hex — imprimable) & le dépôt
porte la clé publique ; une graine perdue attend la rotation de clé (re-signer tous
les arbres, même commit) \\
\bottomrule
\end{tabularx}
\end{table}

\subsection{Le sachet — l'objet physique qui vaut le coffre}

Le sachet est une pochette physique, \textbf{individuelle} (un média du piège par
coffre — borne consignée), qui porte : les \textbf{pièges d'époque}
(\code{piege-CF3-epoque-N.txt}, un fichier par époque, ajoutés à chaque clôture) ; les
\textbf{copies hors-ligne des fragments CF2} (T = 3 : la relève si les gardes machine
brûlent) ; et la \textbf{carte papier}. Il se branche pendant les actes qui l'exigent —
et \emph{seulement} pendant eux : le geste de débrancher est consigné au journal de
cérémonie, parce qu'un média resté branché pendant la session est un média exposé.

\subsection{La carte du sachet — le témoin indécevable}

La carte est une fiche papier, écrite \emph{à la main}, une ligne par clôture :

\begin{lstlisting}
TRIADE - coffre CF3 (carte du sachet)
2026-10-30  epoque 1   checksum 0cc89a96
                                (migration CF2 -> CF3, VERIFIER VERT x2)
<ligne suivante a la prochaine cloture - jamais d'effacement>
\end{lstlisting}

Sa raison d'être : l'anti-retour-en-arrière (§ \ref{sec:temoins}). Un attaquant admin
peut effacer le journal local et vous resservir un conteneur ancien avec son piège
ancien — il ne peut pas imiter votre écriture dans votre sachet. Le jour où
\code{verifier} affiche un checksum absent de votre carte, le mensonge est nommé.
La carte ne protège que si elle est \emph{regardée} : le verdict à froid ré-exécuté
par la machine est le filet (recommandation : machine froide, le lendemain d'une
cérémonie).

\subsection{La frontière des supports — git, \code{.keys}, sachet, Drive}

\begin{table}[htbp]
\centering\small
\caption{La frontière des supports — ce qui vit où, et pourquoi.}
\label{tab:frontiere}
\begin{tabularx}{\textwidth}{@{}l Y Y@{}}
\toprule
\textbf{Support} & \textbf{Ce qu'il porte} & \textbf{Ce qu'il ne porte JAMAIS} \\
\midrule
Dépôt git (poussé, PUBLIC depuis D220) & le code, la spécification signée
(\code{v06/I/}), les preuves publiques du sceau, le journal de travail & aucun
conteneur, aucun fragment, aucune graine, AUCUNE PHRASE (le séquestre 5/3
hors dépôt, D223) : \code{/.keys/} est gitignoré \\
\code{.keys/} du clone & conteneurs CF1/CF2/CF3, gardes CF2, config, journaux et PV
locaux, miroir d'époque & jamais dans git — la ligne existe depuis D12/D27, le
coffre-fort depuis D75 \\
Le sachet & pièges d'époque, T copies hors-ligne, la carte papier & jamais branché
pendant une session de travail \\
Le Drive & le miroir complet du dépôt, le kit de cérémonie (\code{v06/} identique au
miroir par construction, D130), le conteneur CF1 de secours & aucune graine de clé I \\
\bottomrule
\end{tabularx}
\end{table}

\begin{terrain}[— rien à pousser après une cérémonie de coffre]
La question du terrain, juste après la migration CF3 : « je dois pousser quelque chose
sur le git maintenant ? » Non — par construction. Tout ce que la migration écrit vit
dans \code{.keys\textbackslash} (ignoré : ligne 28 du \code{.gitignore}), le piège vit
sur le média, la boîte noire est un \code{*.log} (ignoré). Un \code{git status} rend
\code{nothing to commit} : le clone est exactement ce que GitHub a. Les seuls
\emph{push} de votre chaîne sont les remises du sceau — et ceux-là, la cérémonie du
sceau les guide. La consignation au \emph{journal Git} (le troisième témoin d'époque)
se fait côté laboratoire, par le commit qui accompagne le rapport D122.
\end{terrain}

\subsection{La rotation — le remède quand un secret a pu fuir}\label{sec:rotation}

La phrase commise en git y demeure à vie (D81) ; le média photographié pendant un acte
sert une époque (borne assumée). Le remède n'est jamais d'espérer : c'est de
\emph{tourner}.

\begin{table}[htbp]
\centering\small
\caption{Les rotations du coffre — quoi tourne, quand, et ce qui ne change pas.}
\label{tab:rotations}
\begin{tabularx}{\textwidth}{@{}l Y Y@{}}
\toprule
\textbf{Rotation} & \textbf{Le geste} & \textbf{Ce qui ne change pas} \\
\midrule
Phrase de contrôle & \code{changer-phrase} (l'ancienne devient inopérante, les
entrées intactes) ; le changement passe par la cérémonie sous séquestre —
trois parts assemblées à froid, la phrase neuve re-découpée en CINQ parts,
JAMAIS écrite au dépôt (le refus nommé la garde) & le coffre, les projets,
le sel (le re-chiffrement prend un sel neuf) \\
Moitiés CF3 & chaque \emph{clôture} : S\_serrure et S\_piège frais, époque N+1 &
K\_user, \code{sel\_argon} — la clé utilisateur est INVARIABLE à travers les
époques \\
Secret maître CF2 & re-découpage Shamir (rotation complète du coffre CF2) & la base
JSON, le seuil \\
Clé I du sceau & remplacer \code{cle-publique-I.txt} ET toutes les signatures dans le
MÊME commit (un état mixte est refusé par \code{verifier}) ; la rotation est consignée
au journal & les arbres, les manifestes (re-signés par la nouvelle clé) \\
\bottomrule
\end{tabularx}
\end{table}

La règle de jugement : une exposition \emph{prouvée} (phrase vue, média photographié)
exige la rotation dans la session qui suit ; une exposition \emph{possible} (doute,
reliquat) se traite au calendrier. Entre les deux, votre degré de sommeil est le
dernier arbitre — c'est à vous que la menace s'adresse.
